% DUBTE: l'original dona el camp en «N m^{-1}» (errata probable per N C^{-1} o V m^{-1}); es transcriu tal com és.
Després de diversos mesuraments es va determinar que hi ha un camp elèctric que envolta la Terra. La magnitud d'aquest camp a la superfície terrestre és d'uns $150\ \mathrm{N\,m^{-1}}$ i està dirigit cap al centre de la Terra.

\begin{apartat}{1}
Quin és el valor de la càrrega elèctrica de la Terra? (Considereu tota la càrrega concentrada en un punt al centre del planeta.)
\end{apartat}

\begin{apartat}{1}
Quants electrons de més ha de tenir una gota d'aigua de $18\ \mu\mathrm{m}$ de radi perquè estigui estacionària, és a dir, perquè no caigui, quan es troba a una altura propera a la superfície terrestre? (Considereu que la gota té forma esfèrica.)
\end{apartat}

\dades{%
  Radi de la Terra, $R_\mathrm{T} = 6,37\times 10^{6}\ \mathrm{m}$.\\
  Càrrega de l'electró, $q_\mathrm{e} = -1,60\times 10^{-19}\ \mathrm{C}$.\\
  Densitat de l'aigua, $\rho_\mathrm{aigua} = 1,00\times 10^{3}\ \mathrm{kg\,m^{-3}}$.\\
  $g = 9,81\ \mathrm{m\,s^{-2}}$.\\[2pt]
  $k = \dfrac{1}{4\pi\varepsilon_0} = 8,99\times 10^{9}\ \mathrm{N\,m^2\,C^{-2}}$.}
